% BEGIN MODEL ARTIFACT: model:classical-propositional-logic
\begin{modelbox}{Model (Classical Propositional Logic)}
\modelid{model:classical-propositional-logic}
\modelkind{model}
\modellayer{propositional-logic}
\begin{model}[Classical Propositional Logic]\label{model:classical-propositional-logic}
\modeldepends{Propositional language; valuation; bridge:free-sigma-algebra-evaluation; bridge:consequence}

\modelnote{
The seed is the model datum itself.
\[
  v:\mathsf P\to\lvert\mathbf 2\rvert,
  \qquad
  \lvert\mathbf 2\rvert=\{0,1\}.
\]
}

\modelbinding{
\[
  \Sigma:=\Omega,
  \qquad
  X:=\mathsf P,
  \qquad
  \text{target}:=\mathbf 2,
  \qquad
  \sigma:=v,
  \qquad
  \widehat v:=\operatorname{eval}(\mathbf 2)(v).
\]
}

\begin{modelbridge}
  \bridgeclause{\[
    v\models\varphi
    \quad:\Longleftrightarrow\quad
    \widehat v(\varphi)=1.
  \]}
\end{modelbridge}

\begin{modelconsequences}
  \modelconsequence{Connective recursion is inherited from
  \emph{Free \(\Sigma\)-Algebra and Evaluation}.}
  \modelconsequence{Semantic consequence, validity, and \(\mathrm{Th}(v)\) are
  inherited from \emph{Consequence}.}
\end{modelconsequences}
\end{model}
\end{modelbox}
% END MODEL ARTIFACT: model:classical-propositional-logic
